
\subsection{Gráficos 2D}

% ___________________________________________________________________________

\begin{frame}\frametitle{Representación gráfica}
\begin{itemize}
\item {} 
Repesentar datos gráficamente en Octave es sencillo e intuitivo.
\item  
  Se pueden usar distintos sistemas gráficos a través de un conjunto de funciones compatibles con Matlab.
\item Verremos un ejemplo en 2D
\end{itemize}

\end{frame}

% % ___________________________________________________________________________

\begin{frame}[containsverbatim]
\frametitle{Plot 2D}

La manera más sencilla de representar datos es mediante la función
\emph{plot}
\begin{description}
\item [Ejemplo:]
Representar una curva en el plano a partir de dos series de
datos:
\begin{Octave}
>> x=linspace(0,20,100); 
>> y=exp(-x./100).*sin(x);
>> plot(x,y);
\end{Octave}
\end{description}
\end{frame}

% ___________________________________________________________________________

\begin{frame}\frametitle{Plot 2D (II)}

El resultado (similar a)...
\begin{figure}[htbp]\begin{center}
%\includegraphics[width=600px]{files/abanico.jpg}
\pgfimage[width=0.8\linewidth]{img/plot-exp-sin.png}
\end{center}\end{figure}

\end{frame}

% ___________________________________________________________________________


\begin{frame}[containsverbatim]
\frametitle{Plot 2D (III)}

Los atributos de las gráficas se introducen con la ventana activa
\begin{description}
\item [Ejemplo:]

\begin{Octave}
>> title('Una funcion cualquiera') 
>> xlabel('Tiempo') 
>> ylabel('Amplitud')
\end{Octave}

\end{description}

\end{frame}


% ___________________________________________________________________________

\hypertarget{plot-iv}{}
\pdfbookmark[0]{Plot 2D (IV)}{plot-iv}
\begin{frame}\frametitle{Plot IV}

El resultado...
\begin{figure}[htbp]\begin{center}

\pgfimage[width=0.8\linewidth]{img/abanico2.png}
\end{center}\end{figure}

\end{frame}

% ___________________________________________________________________________

\hypertarget{plot-v}{}
\pdfbookmark[0]{Plot 2D (V)}{plot-v}
\begin{frame}[containsverbatim]
\frametitle{Plot V}
\small
Podemos poner varias curvas con distintos estilos:
\begin{Octave}
>> set(0, 'DefaultLineLineWidth', 1.5);   % Grosor 
>> x = linspace(0,20,50);
>> plot (x,exp(x/10).*sin(x),':',...
         x,exp(x/10).*cos(x),'-^'); 
\end{Octave}

Añadir una leyenda a posteriori...
\begin{Octave}
>> legend('exp(x/10)*sin(x)','exp(x/10)*cos(x)');
>> replot;  % Volver a pintar la gr'afica
\end{Octave}

Y exportar a un fichero
\begin{Octave}
  print -deps dibujo.eps
\end{Octave}

En la ayuda de \verb|print| está la lista de dispositivos (\verb|-d|)
\end{frame}

% ___________________________________________________________________________

\begin{frame}\frametitle{Plot 2D (VI)}
  El resultado...
\begin{center}
\pgfimage[width=0.85\linewidth]{img/plot-exp-sin-cos.png}
\end{center}
\end{frame}

\subsection{Gráficos en 3D}


\begin{frame}[containsverbatim]
  \frametitle{Curvas 3D}
  Las gráficas 3D más sencillas\footnote{Aunque
  \texttt{plot3} tiene otras posibilidades, por ejemplo representar
  familias de curvas (cuando \texttt{x}, \texttt{y}, \texttt{z} son
  matrices). Véase la ayuda.} son curvas trazadas mediante la función
  \emph{plot3}
  \begin{description}
  \item [Ejemplo:] Representar una curva en el espacio a partir de tres
    vectores de datos, \verb|x|, \verb|y|, \verb|z|.
\begin{Octave}
  >> t=0:0.1:10*pi;
  >> x=(1+sin(t)).*cos(t);
  >> y=(1+sin(t)).*sin(t);
  >> z=t;
  >> plot3(x,y,z);
\end{Octave}
\end{description}
\end{frame}

\begin{frame}\frametitle{Curvas 3D: El resultado...}
\begin{center}
\pgfimage[width=0.9\linewidth]{img/plot3.png}
\end{center}
\end{frame}

\begin{frame}[containsverbatim]
  \frametitle{Superficies 3D: Generación de mallas}
  La función \verb|[x,y]=meshgrid(xmin:hx:xmax, ymin:hy:ymax)| genera
  dos matrices \verb|x| e \verb|y|:
  \begin{itemize}
  \item Filas de \verb|x|: copias del vector \verb|xmin:hx:xmax|
  \item Columnas de \verb|y|: vector \verb|(ymin:hy:ymax)'|
  \end{itemize}
  Ejemplo
\begin{Octave}
  >> [x,y]=meshgrid(0:2,0:1)
  x =
   0   1   2
   0   1   2

 y =
   0   0   0
   1   1   1
\end{Octave}
Los \verb|(x[i],y[j])| son un mallado de \verb|[xmin,xmax]|$\times$\verb|[ymin,ymax]|
\end{frame}


\begin{frame}[containsverbatim]
  \frametitle{Dibujo de superficies 3D}
  La orden \verb|mesh(x,y,z)| representa una superficie 3D, siendo
  \begin{itemize}
  \item \verb|x|, \verb|y| dos matrices que representan una malla
  \item \verb|z| otra matriz que representa los valores $f(x,y)$
  \end{itemize}
  Ejemplo: gráfica de la función $f(x,y)=sin(x)*cos(y)$
  \begin{Octave}
    >> [x,y]=meshgrid(-pi:0.1:pi, -pi:0.1:pi);
    >> z=sin(x).*sin(y);
    >> mesh(x,y,z);
  \end{Octave}
  Podemos cambiar el punto de vista con la función \verb|view(a,e)|,
  donde 
  \begin{itemize}
  \item \verb|a| es el \emph{azimut} (ángulo de rotación en el plano xy)
  \item \verb|e| es la \emph{elevación} vertical
  \end{itemize}
  \begin{Octave}
  >> view(30,60);
  >> colorbar('East'); % Pintar una barra vertical
  \end{Octave}
\end{frame}

\begin{frame}\frametitle{Superficies 3D: El resultado...}
\begin{center}
\pgfimage[width=0.9\linewidth]{img/mesh-3d.png}
\end{center}
\end{frame}

\begin{frame}[containsverbatim]
  \frametitle{Curvas de nivel}
  \begin{itemize}
  \item
    La orden \verb|contour(x,y,z)| representa curvas de nivel
  \item \verb|x|, \verb|y| y \verb|z| tienen el mismo significado que
    en la orden \verb|mesh|
  \item Acepta un cuarto parámetro, \verb|n| (número de curvas)
  \item La variante \verb|contourf| dibuja un degradado de colores
  \end{itemize}
  \begin{Octave}
  >> [x,y]=meshgrid(-pi:0.1:pi, -pi:0.1:pi);
  >> z=sin(x).*sin(y);
  >> contourf(x,y,z);    
  \end{Octave}
\end{frame}

\begin{frame}\frametitle{Curvas de nivel: El resultado...}
\begin{center}
\pgfimage[width=0.9\linewidth]{img/contour.png}
\end{center}
\end{frame}


% \begin{frame}[containsverbatim]
% \frametitle{Mallas estructuradas 3D con VTK}
% \begin{enumerate}
% \item Creamos un array 3D y lo grabamos en formato \verb|hdf5|
%   \begin{Octave}
% nx=ny=nz=4; % Por ejemplo, 4 ptos en cada eje
% A=rand(nx,ny,nz); % array aleatorio nx*ny*nz
% save -hdf5 "malla_estruc.h5" A;
%   \end{Octave}
% \item Convertimos el fichero \verb|hdf5| en formato \verb|vtk|
%   \begin{sh}
% $ h5tovtk -d /A/value malla_estruc.h5
%   \end{sh}  % $fin de modo matemático, editor!
%   \begin{itemize}
%   \item \footnotesize Se ha usado el programa \texttt{h5tovtk} (en
%     Ubuntu, paquete h5utils).
%   \item \footnotesize La opción ``\texttt{-d}'' indica la ruta hacia
%     el ``dataset'', dentro del fichero (ésta se puede localizar con
%     el programa \texttt{h5dump} o \texttt{h5ls -r})
%   \end{itemize}
% \item Visualizamos el fichero resultante (\texttt{malla\_estruc.vtk})
%   con \texttt{mayavi2} o con \texttt{paraview} (filtro ``\emph{Treshold}'')
% \end{enumerate}
% \end{frame}
%
% \begin{frame}\frametitle{Resultado (generado con paraview)}
% \begin{center}
% \pgfimage[width=0.9\linewidth]{img/cubo.png}
% \end{center}
% \end{frame}


%%% Local Variables: 
%%% mode: latex
%%% TeX-master: "modulo_octave"
%%% End: 
